\section{The canonical Transformer block}\label{sec:full-block}

The preceding sections reconstructed the two sides of the architectural factorization separately. Attention constructs a receiver interface from the current token state. The residual, normalization, and FFN form a structured receiver-local continuation. This section states the whole block first at that coarse grain, then exposes a finer two-sublayer presentation.

\subsection{Canonical deterministic PreNorm block}

Let
\[
H=(h_i)_{i\in S}\in V^S
\]
be the current token state. Let $N_1,N_2:V\to V$ be receiver-local normalization maps and let $m\in\mathcal H$ index heads. Let $Q_m,K_m,V_m$, and $O_m$ be the declared linear head maps of compatible types. Define
\[
q_i^{(m)}=Q_mN_1(h_i),
\qquad
k_j^{(m)}=K_mN_1(h_j),
\qquad
v_j^{(m)}=V_mN_1(h_j).
\]
Let $\alpha_{ij}^{(m)}(H)$ be masked row-normalized attention coefficients. With optional attention bias $b_{\mathrm{att}}$, define
\begin{equation}\label{eq:token-route}
t_i(H)
:=
b_{\mathrm{att}}
+
\sum_{m\in\mathcal H}O_m\left(\sum_{j\in S}\alpha_{ij}^{(m)}(H)v_j^{(m)}\right).
\end{equation}
The first residual state is
\begin{equation}\label{eq:first-residual}
u_i:=h_i+t_i(H).
\end{equation}
For a two-layer local sublayer,
\begin{equation}\label{eq:local-ffn-block}
f_i(u_i):=W_2\phi(W_1N_2(u_i)+b_1)+b_2,
\end{equation}
and the block output is
\begin{equation}\label{eq:block-output}
h_i^+:=u_i+f_i(u_i).
\end{equation}

\subsection{Coarse receiver factorization}

\begin{theorem}[Receiver-interface and local-continuation decomposition]\label{thm:block-receiver-factorization}
For every receiver $i$, let
\[
Z_i^{\mathrm{att}}
:=
\im\bigl(H\mapsto(h_i,t_i(H))\bigr),
\]
and define the attention-constructed interface, corestricted to its realized image,
\begin{equation}\label{eq:block-coarse-interface}
Q_i^{\mathrm{att}}:V^S\twoheadrightarrow Z_i^{\mathrm{att}},
\qquad
Q_i^{\mathrm{att}}(H):=\bigl(h_i,t_i(H)\bigr).
\end{equation}
Define the local continuation on $Z_i^{\mathrm{att}}$ by
\begin{equation}\label{eq:block-coarse-continuation}
G_i^{\mathrm{block}}(h,t)
:=\nu+\FFN(N_2(\nu)),
\qquad
\nu:=h+t,
\end{equation}
where $\FFN(x)=W_2\phi(W_1x+b_1)+b_2$. Then the receiver branch of the canonical PreNorm block factors as
\[
\boxed{
B_i^{\mathrm{block}}
=
G_i^{\mathrm{block}}Q_i^{\mathrm{att}}.
}
\]
Moreover, $Q_i^{\mathrm{att}}$ has a source-resolved refinement through the head--source contributions in \eqref{eq:token-route}, while $G_i^{\mathrm{block}}$ has the marked hidden-coordinate refinement of \Cref{prop:ffn-hidden-carrier}.
\end{theorem}

\begin{proof}
Evaluating \eqref{eq:block-coarse-continuation} on \eqref{eq:block-coarse-interface} gives $\nu=h_i+t_i(H)=u_i$ and therefore
\[
G_i^{\mathrm{block}}Q_i^{\mathrm{att}}(H)
=u_i+W_2\phi(W_1N_2(u_i)+b_1)+b_2
=h_i^+.
\]
Linearity of each $O_m$ expands the transported exposure as the source-resolved family
\[
t_i(H)
=b_{\mathrm{att}}
+
\sum_{m,j}\alpha_{ij}^{(m)}(H)O_mv_j^{(m)}.
\]
The FFN refinement follows from the marked intermediate activation vector.
\end{proof}

The theorem gives the main architectural reading of the block:
\[
\boxed{
\text{attention constructs }Q_i
\qquad\text{and}\qquad
\text{the residual--FFN subprogram realizes }G_i.
}
\]
Biases remain affine terms. They need not be converted into artificial ``null sources'' merely to force every term into one routed sum.

\subsection{Fine two-sublayer presentation}

The coarse factorization treats the complete block as one receiver branch. A finer presentation marks the intermediate residual state $U=(u_i)_i$ as a successor state of the attention sublayer and predecessor state of the FFN sublayer.

\begin{corollary}[Fine two-stage receiver presentation]\label{cor:fine-two-stage-block}
The canonical block admits the finer composition
\[
H\xrightarrow{B^{\mathrm{att}}}U\xrightarrow{B^{\mathrm{ffn}}}H^+,
\]
where each displayed interface is corestricted to its realized image and
\[
Q_i^{\mathrm{tok}}(H)=\bigl(h_i,t_i(H)\bigr),
\qquad
G_i^{\mathrm{tok}}(h,t)=h+t,
\]
and
\[
Q_i^{\mathrm{ffn}}(U)=\bigl(u_i,f_i(u_i)\bigr),
\qquad
G_i^{\mathrm{res},2}(u,\delta)=u+\delta.
\]
The transported token contribution is source-resolved over head--source pairs. The second residual interface contains the completed local FFN contribution, while that contribution is internally refined by the actual hidden activation
\[
u_i
\longmapsto
a_i=\phi(W_1N_2(u_i)+b_1)
\longmapsto
f_i(u_i)=W_2a_i+b_2.
\]
Thus $Q_i^{\mathrm{ffn}}$ is not itself called a hidden interface: the hidden carrier occurs inside the construction of its second component.
\end{corollary}

\begin{proof}
The first factorization is \eqref{eq:first-residual}. The second is \eqref{eq:block-output}. Their composition equals the coarse branch in \Cref{thm:block-receiver-factorization}.
\end{proof}

The coarse and fine presentations answer different architectural questions. The coarse view identifies the information available to the final receiver continuation. The fine view exposes the intermediate state on which the second local sublayer operates. Neither grain is forced by the external function alone; the selected marks determine which factorization is being claimed.

\begin{corollary}[Gated local continuation]\label{cor:gated-block}
If \eqref{eq:local-ffn-block} is replaced by the gated map \eqref{eq:gated-ffn}, the coarse and fine block factorizations remain valid with hidden coefficients
\[
a_e(u_i)
=
\phi((W_gN_2(u_i)+b_g)_e)(W_uN_2(u_i)+b_u)_e.
\]
\end{corollary}

\subsection{Depth as repeated interface construction}

Let $B_\ell$ be the $\ell$th block and
\[
H^{(\ell+1)}=B_\ell(H^{(\ell)}).
\]
At each layer, the current state determines the evaluated interface constructors
\[
H^{(\ell)}
\longmapsto
\bigl(Q_{\ell,i}^{\mathrm{att}}(H^{(\ell)})\bigr)_i
\]
and, in the additive regime, a token-routing operator family
\[
H^{(\ell)}\longmapsto A_{\ell,H^{(\ell)}}.
\]
In the positive scalar specialization it also determines receiver- and head-relative potentials
\[
H^{(\ell)}\longmapsto\mathcal U_{\ell,H^{(\ell)}}.
\]
Ordinary function composition is enough to describe depth:
\[
H^{(0)}
\xrightarrow{B_0}
H^{(1)}
\xrightarrow{B_1}
\cdots
\xrightarrow{B_{L-1}}
H^{(L)}.
\]
No category of process cuts is required for the main reconstruction. Whether these represented transitions descend coherently from an implementation or remain equivalent under a change of cut is a separate question not assumed here.

\subsection{Conditionally variational routing, globally driven computation}

Within one positive attention row, the allocation $\alpha_{ri}^{(m)}$ uniquely minimizes the receiver free-energy functional of \Cref{thm:receiver-free-energy}. The whole block does not minimize one fixed global free energy. Its update changes the represented state used to construct the next scores, masks, values, interfaces, and local hidden coefficients.

The variational statement is therefore receiver- and state-conditional: each declared positive row minimizes its own current free-energy functional, while the block update changes the state from which later rows are constructed. No energy-conservation law, monotone global free-energy descent, or physical equilibrium interpretation follows.

\subsection{Residual continuation and schedule variants}

The residual path marks direct availability of the predecessor-local state:
\[
h_i\longmapsto h_i+t_i\longmapsto h_i+t_i+f_i.
\]
It prevents the routed aggregate from becoming the receiver's sole successor variable. This does not guarantee invertibility or information conservation; distinct predecessor states may still share one successor.

The exact $Q/G$ split depends on schedule. In a PostNorm block, normalization occurs after residual addition. In a parallel block, attention and FFN branches may receive the same predecessor state rather than a sequentially updated one. Reversed, recurrent, adaptive-depth, asynchronous, or carrier-changing schedules give different receiver-wise factorizations because they change which represented state is supplied to each interface constructor and continuation. The architecture is the marked schedule of these factorizations, not merely the multiset of formulas appearing in the block.
